\section{Conclusiones}\label{sec:conclusiones}

El modelo final es deliberadamente simple. Usa el régimen de la etiqueta, que explica la mayor parte del error que se
puede reducir, y el rendimiento de los vecinos del mismo régimen a 400\,m, combinados con un proceso gaussiano. Cada
pieza entró porque mejoró la validación de forma consistente y pasó una prueba de permutación, y cada pieza que se
quedó fuera, como las variables satelitales y climáticas o los modelos de árboles, lo hizo porque no mejoró. El
modelo obtiene 0.458\,t/ha en validación cruzada, y estimamos un RMSE de 0.46 a 0.49\,t/ha en las 59 parcelas, con un
margen adicional por el tamaño de la muestra de calificación.

La lección principal es que en este reto el límite no lo pone el algoritmo sino las etiquetas. Con dos formas de
reportar y rendimientos declarados, la mejora sobre la media por régimen es pequeña y cualquier variable adicional
compite contra el ruido. Mejores datos de campo, como el rendimiento medido o la fecha de siembra, ayudarían más que
un modelo más complejo.
